% =========================================================
% Volume IV  \S{}1.7.?  --- C as a Model of the Field Signature
% WIRING NOTE (cross-volume):
%   Verify exact labels before final commit.
%
% Expected targets:
%   def:arithmetic-signature
%   def:algebraic-field
%   thm:c-is-field
%   thm:fundamental-theorem-of-algebra
%   thm:r-embeds-in-c
%   def:structure-expansion
%   def:structure-reduction
% =========================================================

\subsection*{The Complex Numbers as a Structure}

\begin{remark*}[Modeling Goal]
We exhibit $\mathbb{C}$ as a concrete field model. The real numbers embed into
$\mathbb{C}$, and the resulting structure is the first standard number system
that is algebraically closed. Unlike $\mathbb{R}$, however, $\mathbb{C}$ admits
no compatible ordered-field expansion.
\end{remark*}

% ---------------------------------------------------------
% Signature
% ---------------------------------------------------------

\begin{remark*}[Exposition]

A model consists of a signature together with an interpretation on a carrier.

The signature considered here is the field signature

\[
\mathcal L_{\mathrm{field}}
=
\{+,\cdot,-,{}^{-1},0,1\}.
\]

It contains the field operations and distinguished constants but no order
relation.

\end{remark*}

% ---------------------------------------------------------
% Model
% ---------------------------------------------------------

\begin{definition}[The Complex Number Model $\mathcal C$]
\label{def:complex-number-model}

The \textbf{complex number model} is the
$\mathcal L_{\mathrm{field}}$-structure

\[
\mathcal C
=
\bigl(
\mathbb C;
+^{\mathcal C},
\cdot^{\mathcal C},
-^{\mathcal C},
{}^{-1\,\mathcal C},
0^{\mathcal C},
1^{\mathcal C}
\bigr),
\]

whose carrier

\[
|\mathcal C|
=
\mathbb C
\]

is the set of complex numbers constructed in Volume~II, and whose symbols
are interpreted by the usual complex-number operations.

\end{definition}

% ---------------------------------------------------------
% Quantified statement
% ---------------------------------------------------------

\begin{remark*}[Standard quantified statement]

\[
|\mathcal C|
=
\mathbb C,
\]

\[
+^{\mathcal C},
\cdot^{\mathcal C}
:
\mathbb C\times\mathbb C
\to
\mathbb C,
\]

\[
-^{\mathcal C}
:
\mathbb C
\to
\mathbb C,
\]

\[
{}^{-1\,\mathcal C}
:
\mathbb C\setminus\{0\}
\to
\mathbb C,
\]

\[
0^{\mathcal C},
1^{\mathcal C}
\in
\mathbb C.
\]

\end{remark*}

% ---------------------------------------------------------
% Dependencies
% ---------------------------------------------------------% ---------------------------------------------------------
% Interpretation
% ---------------------------------------------------------

\begin{remark*}[Interpretation]

The structure

\[
\mathcal C
\]

packages the complex-number system into a single mathematical object.

The carrier and operations are those already constructed in Volume~II. The
model merely records how the field signature is interpreted on that carrier.

\end{remark*}

% ---------------------------------------------------------
% Satisfaction
% ---------------------------------------------------------

\begin{remark*}[Satisfaction]

The structure

\[
\mathcal C
\]

is a model of the field axioms:

\[
\mathcal C
\models
\Sigma_{\mathrm{field}},
\]

equivalently,

\[
\operatorname{Field}
(\mathbb C,+,\cdot,0,1).
\]

This is witnessed by

\hyperref[thm:c-is-field]
{$\mathbb C$ Is a Field}

in Volume~II.

\end{remark*}

% ---------------------------------------------------------
% Algebraic closure
% ---------------------------------------------------------

\begin{remark*}[Algebraic Closure]

The most important structural property of the complex numbers is not merely
that they form a field.

The defining feature of $\mathbb C$ is that it is algebraically closed.

Equivalently,

\[
\forall p(x)\in\mathbb C[x],
\qquad
\deg p \ge 1
\Longrightarrow
\exists z\in\mathbb C
\quad
p(z)=0.
\]

This property is witnessed by

\hyperref[thm:fundamental-theorem-of-algebra]
{Fundamental Theorem of Algebra}.

\end{remark*}

\begin{remark*}[Interpretation]

Every nonconstant polynomial with complex coefficients has a complex root.

Consequently every polynomial factors completely into linear factors over
$\mathbb C$.

This is the structural reason the complex numbers occupy the final stage of the
standard algebraic number-system hierarchy.

\end{remark*}

% ---------------------------------------------------------
% Real embedding
% ---------------------------------------------------------

\begin{remark*}[Real Numbers as a Substructure]

The real numbers embed into the complex numbers through the canonical map

\[
r
\longmapsto
r+0i.
\]

This embedding is witnessed by

\hyperref[thm:r-embeds-in-c]
{$\mathbb R$ Embeds in $\mathbb C$}.

Under this embedding the real numbers appear as a subfield of the complex
field.

\end{remark*}

% ---------------------------------------------------------
% Reducts
% ---------------------------------------------------------

\begin{remark*}[Reduction to the Additive Group]

Forgetting multiplication leaves

\[
(\mathbb C;+,-,0),
\]

the additive abelian group of complex numbers.

\end{remark*}

\begin{remark*}[Reduction to the Multiplicative Group]

Restricting to nonzero elements and forgetting addition leaves

\[
(\mathbb C^\times;\cdot,1),
\]

the multiplicative abelian group of nonzero complex numbers.

\end{remark*}

% ---------------------------------------------------------
% Ordered field impossibility
% ---------------------------------------------------------

\begin{remark*}[No Ordered-Field Expansion]

Unlike $\mathbb Q$ and $\mathbb R$, the complex field cannot be expanded into
an ordered field compatible with its field operations.

Indeed,

\[
i^2=-1.
\]

In every ordered field, every square is nonnegative. Hence

\[
i^2 \ge 0.
\]

But this would imply

\[
-1 \ge 0,
\]

which contradicts the ordered-field axioms.

Therefore no order relation can be added to the complex field so as to make it
an ordered field.

\end{remark*}

% ---------------------------------------------------------
% Blueprint position
% ---------------------------------------------------------

\begin{remark*}[Blueprint Position]

The complex-number model occupies the final stage of the standard number
hierarchy:

\[
\mathbb N
\hookrightarrow
\mathbb Z
\hookrightarrow
\mathbb Q
\hookrightarrow
\mathbb R
\hookrightarrow
\mathbb C.
\]

The passage

\[
\mathbb R
\longrightarrow
\mathbb C
\]

supplies solutions to polynomial equations that have no real solutions.

For example,

\[
x^2+1=0
\]

has no solution in $\mathbb R$ but has the solutions

\[
x=\pm i
\]

in $\mathbb C$.

\end{remark*}

% ---------------------------------------------------------
% Structural significance
% ---------------------------------------------------------

\begin{remark*}[Structural Significance]

The complex-number model is

\begin{itemize}
    \item a field,
    \item an algebraically closed field,
    \item not an ordered field.
\end{itemize}

The real numbers complete the ordered-field hierarchy.

The complex numbers complete the algebraic hierarchy.

For this reason $\mathbb C$ serves as the canonical ambient field for algebra,
complex analysis, algebraic geometry, and much of modern mathematics.

\end{remark*}

\begin{dependencies}

\begin{itemize}
    \item \hyperref[def:algebraic-field]{Field}.
    \item \hyperref[def:complex-numbers]{The Complex Numbers}.
    \item \hyperref[def:addition-on-c-classes]{Addition on $\mathbb C$}.
    \item \hyperref[def:multiplication-on-c-classes]{Multiplication on $\mathbb C$}.
    \item \hyperref[thm:c-is-a-field]{$\mathbb C$ Is a Field}.
    \item TODO: formalize the Fundamental Theorem of Algebra.
\end{itemize}

\end{dependencies}
